% GENERATED FILE. Edit canonical lessons, then run npm run generate:books.
% canonical-source-hash: fnv1a64:1dad4ab40bc77a23
% canonical-lessons: TE-C513-vyadhi, TE-C513-antuvyadhi, TE-C513-gundepotu, TE-C513-paksavatam, TE-C513-ubbasam

\chapter{Diseases}
\label{ch:te-diseases}
\hlchaptermodality{513}

\begin{chapteropening}
\textbf{By the end of this chapter:} \emph{I can name some serious and contagious diseases.}

Complete the last lesson of chapter 513: I can name some serious and contagious diseases.
\end{chapteropening}

\section[\texorpdfstring{\te{వ్యాధి} (vyādhi)}{వ్యాధి (vyādhi)}]{\te{వ్యాధి} (vyādhi) — a disorder; a disease}

\label{lesson:TE-C513-vyadhi}



\begin{quote}
Before the new one: say the Telugu for to be bedridden, then the Telugu for to come round.
\end{quote}

\subsection*{You'll want to know: \te{వ్యాధి}}
\textbf{\te{వ్యాధి}} — \emph{vyādhi} — \textquotedblleft{}a disorder; a disease\textquotedblright{}.

\begin{quote}
\te{ఆ} \te{వ్యాధికి} \te{ఇంకా} \te{మందు} \te{లేదు}.
\end{quote}

\begin{quote}
\emph{ā vyādhiki iṅkā mandu lēdu.}
\end{quote}

\begin{quote}
There is still no cure for that disease.
\end{quote}

More formal than \te{రోగం} or \te{జబ్బు}; doctors and the news say \te{వ్యాధి}.

\begin{cousinweb}
\textbf{\te{వ్యాధి}} (\emph{vyādhi}) is borrowed from Sanskrit \emph{vyādhi}, \textquotedblleft{}disease\textquotedblright{} (CDIAL 12202).
\end{cousinweb}

\begin{grammarlens}[title={What you've built}]
One more word for describing what you lived and felt.
\end{grammarlens}

\subsection*{Your turn}
\begin{itemize}
\raggedright
  \item \emph{Say it:} \emph{vyādhi}
  \item \emph{Say it:} \emph{vyādhi}, once more
  \item \emph{Say it:} \emph{ā vyādhiki iṅkā mandu lēdu.} — There is still no cure for that disease.
  \item \emph{Recall:} say the Telugu for to be bedridden, then the Telugu for to come round, then say \emph{vyādhi} again
\end{itemize}

\begin{tcolorbox}[breakable,colback=teal!4,colframe=teal!35!black,title={Before you move on}]
What does \te{వ్యాధి} mean? (\textquotedblleft{}A disorder; a disease\textquotedblright{}.) Say it once more.
\end{tcolorbox}
\section[\texorpdfstring{\te{అంటువ్యాధి} (aṇṭuvyādhi)}{అంటువ్యాధి (aṇṭuvyādhi)}]{\te{అంటువ్యాధి} (aṇṭuvyādhi) — a contagious disease}

\label{lesson:TE-C513-antuvyadhi}



\begin{quote}
Before the new one: say the Telugu for to come round, then the Telugu for a disorder.
\end{quote}

\subsection*{You'll want to know: \te{అంటువ్యాధి}}
\textbf{\te{అంటువ్యాధి}} — \emph{aṇṭuvyādhi} — \textquotedblleft{}a contagious disease\textquotedblright{}.

\begin{quote}
\te{అంటువ్యాధి} \te{వచ్చినప్పుడు} \te{బడులు} \te{మూసేశారు}.
\end{quote}

\begin{quote}
\emph{aṇṭuvyādhi vaccinappuḍu baḍulu mūsēśāru.}
\end{quote}

\begin{quote}
When the contagious disease broke out, the schools were shut.
\end{quote}

\te{అంటు}, contact, + \te{వ్యాధి}: a disease that spreads by contact.

\begin{cousinweb}
\textbf{\te{అంటువ్యాధి}} (\emph{aṇṭuvyādhi}) joins native \emph{aṇṭu}, \textquotedblleft{}to touch; contact, defilement by touch\textquotedblright{} (DEDR 120), to Sanskrit \emph{vyādhi}, \textquotedblleft{}disease\textquotedblright{} (CDIAL 12202).
\end{cousinweb}

\begin{grammarlens}[title={What you've built}]
One more word for describing what you lived and felt.
\end{grammarlens}

\subsection*{Your turn}
\begin{itemize}
\raggedright
  \item \emph{Say it:} \emph{aṇṭuvyādhi}
  \item \emph{Say it:} \emph{aṇṭuvyādhi}, once more
  \item \emph{Say it:} \emph{aṇṭuvyādhi vaccinappuḍu baḍulu mūsēśāru.} — When the contagious disease broke out, the schools were shut.
  \item \emph{Recall:} say the Telugu for to come round, then the Telugu for a disorder, then say \emph{aṇṭuvyādhi} again
\end{itemize}

\begin{tcolorbox}[breakable,colback=teal!4,colframe=teal!35!black,title={Before you move on}]
What does \te{అంటువ్యాధి} mean? (\textquotedblleft{}A contagious disease\textquotedblright{}.) Say it once more.
\end{tcolorbox}
\section[\texorpdfstring{\te{గుండెపోటు} (guṇḍepōṭu)}{గుండెపోటు (guṇḍepōṭu)}]{\te{గుండెపోటు} (guṇḍepōṭu) — a heart attack}

\label{lesson:TE-C513-gundepotu}



\begin{quote}
Before the new one: say the Telugu for a disorder, then the Telugu for a contagious disease.
\end{quote}

\subsection*{You'll want to know: \te{గుండెపోటు}}
\textbf{\te{గుండెపోటు}} — \emph{guṇḍepōṭu} — \textquotedblleft{}a heart attack\textquotedblright{}.

\begin{quote}
\te{మా} \te{తాతయ్యకు} \te{గత} \te{సంవత్సరం} \te{గుండెపోటు} \te{వచ్చింది}.
\end{quote}

\begin{quote}
\emph{mā tātayyaku gata saṁvatsaraṁ guṇḍepōṭu vaccindi.}
\end{quote}

\begin{quote}
My grandfather had a heart attack last year.
\end{quote}

A heart attack comes: \te{గుండెపోటు} \te{వచ్చింది}. \te{పోటు} is a stab or a throb.

\begin{cousinweb}
\textbf{\te{గుండెపోటు}} (\emph{guṇḍepōṭu}) is native Dravidian: DEDR 1693 lists Telugu \emph{guṇḍe}, \textquotedblleft{}the chest, the heart, courage\textquotedblright{}, with Tamil \emph{kuṇṭi}, \textquotedblleft{}heart\textquotedblright{}; \emph{pōṭu} is a stab of pain.
\end{cousinweb}

\begin{grammarlens}[title={What you've built}]
One more word for describing what you lived and felt.
\end{grammarlens}

\subsection*{Your turn}
\begin{itemize}
\raggedright
  \item \emph{Say it:} \emph{guṇḍepōṭu}
  \item \emph{Say it:} \emph{guṇḍepōṭu}, once more
  \item \emph{Say it:} \emph{mā tātayyaku gata saṁvatsaraṁ guṇḍepōṭu vaccindi.} — My grandfather had a heart attack last year.
  \item \emph{Recall:} say the Telugu for a disorder, then the Telugu for a contagious disease, then say \emph{guṇḍepōṭu} again
\end{itemize}

\begin{tcolorbox}[breakable,colback=teal!4,colframe=teal!35!black,title={Before you move on}]
What does \te{గుండెపోటు} mean? (\textquotedblleft{}A heart attack\textquotedblright{}.) Say it once more.
\end{tcolorbox}
\section[\texorpdfstring{\te{పక్షవాతం} (pakṣavātaṁ)}{పక్షవాతం (pakṣavātaṁ)}]{\te{పక్షవాతం} (pakṣavātaṁ) — paralysis; a stroke}

\label{lesson:TE-C513-paksavatam}



\begin{quote}
Before the new one: say the Telugu for a contagious disease, then the Telugu for a heart attack.
\end{quote}

\subsection*{You'll want to know: \te{పక్షవాతం}}
\textbf{\te{పక్షవాతం}} — \emph{pakṣavātaṁ} — \textquotedblleft{}paralysis; a stroke\textquotedblright{}.

\begin{quote}
\te{పక్షవాతం} \te{వల్ల} \te{ఆయన} \te{కుడి} \te{చెయ్యి} \te{పనిచేయదు}.
\end{quote}

\begin{quote}
\emph{pakṣavātaṁ valla āyana kuḍi ceyyi panicēyadu.}
\end{quote}

\begin{quote}
Because of the stroke, his right hand does not work.
\end{quote}

\te{పక్షం}, side, + \te{వాతం}, wind: in old medicine, a wind that strikes one side.

\begin{cousinweb}
\textbf{\te{పక్షవాతం}} (\emph{pakṣavātaṁ}) is a learned compound built from Sanskrit \emph{pakṣa}, \textquotedblleft{}wing, side\textquotedblright{} (CDIAL 7627), and \emph{vāta}, \textquotedblleft{}wind; wind as a humour of the body\textquotedblright{} (CDIAL 11491).
\end{cousinweb}

\begin{grammarlens}[title={What you've built}]
One more word for describing what you lived and felt.
\end{grammarlens}

\subsection*{Your turn}
\begin{itemize}
\raggedright
  \item \emph{Say it:} \emph{pakṣavātaṁ}
  \item \emph{Say it:} \emph{pakṣavātaṁ}, once more
  \item \emph{Say it:} \emph{pakṣavātaṁ valla āyana kuḍi ceyyi panicēyadu.} — Because of the stroke, his right hand does not work.
  \item \emph{Recall:} say the Telugu for a contagious disease, then the Telugu for a heart attack, then say \emph{pakṣavātaṁ} again
\end{itemize}

\begin{tcolorbox}[breakable,colback=teal!4,colframe=teal!35!black,title={Before you move on}]
What does \te{పక్షవాతం} mean? (\textquotedblleft{}Paralysis; a stroke\textquotedblright{}.) Say it once more.
\end{tcolorbox}
\section[\texorpdfstring{\te{ఉబ్బసం} (ubbasaṁ)}{ఉబ్బసం (ubbasaṁ)}]{\te{ఉబ్బసం} (ubbasaṁ) — asthma}

\label{lesson:TE-C513-ubbasam}



\begin{quote}
Before the new one: say the Telugu for a heart attack, then the Telugu for paralysis.
\end{quote}

\subsection*{You'll want to know: \te{ఉబ్బసం}}
\textbf{\te{ఉబ్బసం}} — \emph{ubbasaṁ} — \textquotedblleft{}asthma\textquotedblright{}.

\begin{quote}
\te{చలికాలంలో} \te{అమ్మకు} \te{ఉబ్బసం} \te{ఎక్కువవుతుంది}.
\end{quote}

\begin{quote}
\emph{calikālaṁlō ammaku ubbasaṁ ekkuvavutundi.}
\end{quote}

\begin{quote}
In winter mum's asthma gets worse.
\end{quote}

Also called \te{ఆస్తమా}, from English.

\begin{cousinweb}
\textbf{\te{ఉబ్బసం}} (\emph{ubbasaṁ}) is native Dravidian: DEDR 632 lists Telugu \emph{ubbasamu}, \textquotedblleft{}asthma\textquotedblright{}, and marks Tamil \emph{uppacam} as borrowed from Telugu.
\end{cousinweb}

\begin{grammarlens}[title={What you've built}]
One more word for describing what you lived and felt.
\end{grammarlens}

\subsection*{Your turn}
\begin{itemize}
\raggedright
  \item \emph{Say it:} \emph{ubbasaṁ}
  \item \emph{Say it:} \emph{ubbasaṁ}, once more
  \item \emph{Say it:} \emph{calikālaṁlō ammaku ubbasaṁ ekkuvavutundi.} — In winter mum's asthma gets worse.
  \item \emph{Recall:} say the Telugu for a heart attack, then the Telugu for paralysis, then say \emph{ubbasaṁ} again
\end{itemize}

\begin{tcolorbox}[breakable,colback=teal!4,colframe=teal!35!black,title={Before you move on}]
What does \te{ఉబ్బసం} mean? (\textquotedblleft{}Asthma\textquotedblright{}.) Say it once more.
\end{tcolorbox}
